\chapter{Algèbre et équations}\label{ch:g9:algebra}

Les lettres désignent des nombres, et calculer avec des lettres prouve
des faits sur \emph{tous} les nombres à la fois. Ce chapitre pratique le
développement et la factorisation --- y compris les trois identités qui
vous accompagneront à travers toutes les mathématiques --- et résout des
équations du premier degré, des équations-produits et des inégalités
simples, toujours une étape justifiée à la fois.

\section{Développer}

\begin{definition}[Développer]\label{def:g9:algebra:expand}
\emph{Développer}\index{développer} signifie transformer un produit en
une somme, en utilisant la distributivité~:
\[
k(a + b) = ka + kb,
\qquad
(a + b)(c + d) = ac + ad + bc + bd .
\]
\end{definition}

\begin{example}\label{ex:g9:algebra:expand}
Chaque terme de la première parenthèse multiplie chaque terme de la
seconde~:
\begin{align*}
(x + 3)(2x + 5)
&= x \times 2x + x \times 5 + 3 \times 2x + 3 \times 5 \\
&= 2x^2 + 5x + 6x + 15 \\
&= 2x^2 + 11x + 15 .
\end{align*}
Les signes voyagent avec les termes~:
\begin{align*}
(3x - 2)(x - 4)
&= 3x^2 - 12x - 2x + 8
= 3x^2 - 14x + 8 .
\end{align*}
\end{example}

\begin{theorem}[Les trois identités]\label{thm:g9:algebra:identities}
Pour tous nombres $a$ et $b$~:
\[
(a+b)^2 = a^2 + 2ab + b^2,
\qquad
(a-b)^2 = a^2 - 2ab + b^2,
\qquad
(a+b)(a-b) = a^2 - b^2 .
\]
\end{theorem}

\begin{proof}
Développer chaque membre de gauche. Pour la première~:
$(a+b)(a+b) = a^2 + ab + ba + b^2 = a^2 + 2ab + b^2$. Pour la seconde,
remplacer $b$ par $-b$ dans la première. Pour la troisième~:
$(a+b)(a-b) = a^2 - ab + ba - b^2 = a^2 - b^2$ (les termes croisés
s'annulent).
\end{proof}

\begin{omfigure}
\begin{tikzpicture}[scale=0.60]
  \draw[thick] (0,0) rectangle (5,5);
  \draw[thick] (3.3,0) -- (3.3,5);
  \draw[thick] (0,3.3) -- (5,3.3);
  \fill[omDef!14] (0,0) rectangle (3.3,3.3);
  \fill[omThm!14] (3.3,3.3) rectangle (5,5);
  \fill[omProp!14] (3.3,0) rectangle (5,3.3);
  \fill[omProp!14] (0,3.3) rectangle (3.3,5);
  \node[font=\small] at (1.65,1.65) {$a^2$};
  \node[font=\small] at (4.15,4.15) {$b^2$};
  \node[font=\small] at (4.15,1.65) {$ab$};
  \node[font=\small] at (1.65,4.15) {$ab$};
  \node[below, font=\small] at (1.65,0) {$a$};
  \node[below, font=\small] at (4.15,0) {$b$};
  \node[left, font=\small] at (0,1.65) {$a$};
  \node[left, font=\small] at (0,4.15) {$b$};
\end{tikzpicture}

{\small L'identité $(a+b)^2 = a^2 + 2ab + b^2$, vue comme aires~: le
grand carré de côté $a+b$ est fait de deux carrés et de deux rectangles
égaux. Oublier le «~$2ab$~», c'est oublier les deux rectangles~!}
\end{omfigure}

\begin{example}\label{ex:g9:algebra:identities}
\[
(x + 5)^2 = x^2 + 10x + 25, \qquad
(3x - 1)^2 = 9x^2 - 6x + 1, \qquad
(2x + 7)(2x - 7) = 4x^2 - 49 .
\]
Calcul mental~: $102^2 = (100 + 2)^2 = 10000 + 400 + 4 = 10404$, et
$98 \times 102 = (100-2)(100+2) = 10000 - 4 = 9996$.
\end{example}

\section{Factoriser}

\begin{method}[Factoriser]\label{met:g9:algebra:factor}
Pour factoriser une expression, essayer dans l'ordre~:
\begin{enumerate}
  \item \emph{facteur commun}~: $ka + kb = k(a + b)$, où $k$ peut
        lui-même être une parenthèse~;
  \item \emph{différence de carrés}~: $a^2 - b^2 = (a+b)(a-b)$~;
  \item \emph{carré parfait}~: $a^2 + 2ab + b^2 = (a+b)^2$ (et avec
        $-2ab$, $(a-b)^2$).
\end{enumerate}
Toujours vérifier en développant le résultat.
\end{method}

\begin{example}\label{ex:g9:algebra:factor}
\emph{Facteur commun~:}
$15x^2 + 10x = 5x \times 3x + 5x \times 2 = 5x(3x + 2)$.

\emph{Parenthèse commune~:}
\[
(x+2)(x-1) + (x+2)(3x+4) = (x+2)\bigl[(x-1) + (3x+4)\bigr]
= (x+2)(4x+3).
\]

\emph{Différence de carrés~:}
$x^2 - 81 = (x+9)(x-9)$, et
\[
(2x+1)^2 - 9 = (2x+1)^2 - 3^2 = (2x+1+3)(2x+1-3) = (2x+4)(2x-2).
\]
Chaque facteur a encore un facteur commun~: $(2x+4)(2x-2) = 2(x+2) \times
2(x-1) = 4(x+2)(x-1)$.
\end{example}

\section{Équations}

\begin{method}[Équations du premier degré]\label{met:g9:algebra:firstdegree}
Pour résoudre une équation comme $5x - 3 = 2x + 9$~:
\begin{enumerate}
  \item rassembler les termes en $x$ d'un côté, les nombres de l'autre,
        en ajoutant la même quantité aux deux membres~;
  \item réduire chaque membre~;
  \item diviser les deux membres par le coefficient de $x$~;
  \item vérifier la solution dans l'équation originale.
\end{enumerate}
\end{method}

\begin{example}\label{ex:g9:algebra:firstdegree}
\begin{align*}
5x - 3 &= 2x + 9 \\
5x - 2x &= 9 + 3 && \text{(ajouter $3 - 2x$ aux deux membres)}\\
3x &= 12 \\
x &= 4 && \text{(diviser par 3).}
\end{align*}
Vérification~: $5 \times 4 - 3 = 17$ et $2 \times 4 + 9 = 17$. La
solution est $4$.

Avec des parenthèses, développer d'abord~: $3(x - 2) = x + 8$ devient
$3x - 6 = x + 8$, donc $2x = 14$ et $x = 7$.
\end{example}

\begin{theorem}[Règle du produit nul]\label{thm:g9:algebra:zeroproduct}
Un produit est nul si et seulement si au moins un facteur est nul~:
\[
A \times B = 0
\quad\Longleftrightarrow\quad
A = 0 \ \text{ ou } \ B = 0 .
\]
\end{theorem}

\begin{proof}
Si un facteur est nul le produit l'est clairement. Réciproquement si
$A B = 0$ avec $A \neq 0$, diviser les deux membres par $A$ donne
$B = 0$.
\end{proof}

\begin{example}\label{ex:g9:algebra:zeroproduct}
Résoudre $(x - 3)(2x + 10) = 0$~: soit $x - 3 = 0$ soit $2x + 10 = 0$,
donc les solutions sont $3$ et $-5$.

Résoudre $x^2 - 49 = 0$~: factoriser d'abord, $(x-7)(x+7) = 0$,
solutions $7$ et $-7$.

Résoudre $x^2 = 3x$~: ne jamais diviser par $x$~! Déplacer et
factoriser~: $x^2 - 3x = x(x - 3) = 0$, solutions $0$ et $3$.
\end{example}

\section{Inégalités}

\begin{proposition}[Règles pour les inégalités]\label{prop:g9:algebra:ineq}
Une inégalité est conservée lorsqu'on ajoute le même nombre aux deux
membres, et lorsqu'on multiplie les deux membres par le même nombre
\emph{positif}. Multiplier les deux membres par un nombre \emph{négatif}
renverse le signe de l'inégalité.
\end{proposition}

\begin{proof}
Si $a \leq b$ alors $b - a \geq 0$. En ajoutant $c$~:
$(b+c) - (a+c) = b - a \geq 0$, donc $a + c \leq b + c$. En multipliant
par $c > 0$~: $bc - ac = (b-a)c \geq 0$, donc $ac \leq bc$. En
multipliant par $c < 0$~: $(b-a)c \leq 0$, donc $ac \geq bc$~: l'ordre
est renversé.
\end{proof}

\begin{example}\label{ex:g9:algebra:ineq}
Résoudre $7 - 2x < 15$~:
\begin{align*}
7 - 2x &< 15 \\
-2x &< 8 && \text{(soustraire 7)} \\
x &> -4 && \text{(diviser par $-2$~: renverser le signe).}
\end{align*}
Les solutions sont tous les nombres strictement plus grands que $-4$~:
sur une droite numérique, un point creux en $-4$ et un hachurage vers la
droite.
\end{example}

\begin{omfigure}
\begin{tikzpicture}[scale=1.0]
  \draw[-Stealth] (-6.4,0) -- (2.6,0) node[below, font=\small] {$x$};
  \foreach \x in {-6,-5,-4,-3,-2,-1,0,1,2}
    \draw (\x,-0.07) -- (\x,0.07);
  \node[below, font=\small] at (-4,-0.12) {$-4$};
  \node[below, font=\small] at (0,-0.12) {$0$};
  \draw[omDef, very thick] (-4,0.02) -- (2.5,0.02);
  \draw[omDef, fill=white, thick] (-4,0.02) circle (2.4pt);
\end{tikzpicture}

{\small Les solutions de $7 - 2x < 15$~: tous les $x > -4$ (point
creux~: $-4$ lui-même n'est pas solution).}
\end{omfigure}

\section{Exercices}

\begin{exercise}[$\star$]\label{exo:g9:algebra:1}
Développer et réduire~:
\[
4(3x - 2), \qquad
(x + 6)(x + 2), \qquad
(2x - 3)(x + 5), \qquad
(x - 1)(x - 9).
\]
\end{exercise}

\begin{exercise}[$\star$]\label{exo:g9:algebra:2}
Développer à l'aide des identités~:
\[
(x + 8)^2, \qquad
(3x - 4)^2, \qquad
(x + 10)(x - 10), \qquad
(5x + 2)(5x - 2).
\]
\end{exercise}

\begin{exercise}[$\star$]\label{exo:g9:algebra:3}
Calculer mentalement, à l'aide d'une identité~: $101^2$~; \quad $59^2$
(écrire $59 = 60 - 1$)~; \quad $53 \times 47$.
\end{exercise}

\begin{exercise}[$\star$]\label{exo:g9:algebra:4}
Factoriser~:
\[
9x + 12, \qquad
x^2 - 64, \qquad
x^2 + 14x + 49, \qquad
18x^2 - 6x .
\]
\end{exercise}

\begin{exercise}[$\star$]\label{exo:g9:algebra:5}
Résoudre, en écrivant chaque étape~:
\[
4x + 7 = 19, \qquad
6x - 5 = 2x + 11, \qquad
5(x - 3) = 3x + 1 .
\]
\end{exercise}

\begin{exercise}[$\star\star$]\label{exo:g9:algebra:6}
Résoudre en utilisant la règle du produit nul~:
\[
(x + 4)(3x - 12) = 0, \qquad
x^2 - 121 = 0, \qquad
(2x - 1)^2 = 0, \qquad
x^2 = 8x .
\]
\end{exercise}

\begin{exercise}[$\star\star$]\label{exo:g9:algebra:7}
Factoriser $E = (3x + 2)^2 - 25$, puis résoudre $E = 0$.
\end{exercise}

\begin{exercise}[$\star\star$]\label{exo:g9:algebra:8}
Résoudre les inégalités et représenter les solutions sur une droite
numérique~:
\[
3x - 4 \geq 8, \qquad
5 - 2x > 1, \qquad
4(x + 1) \leq x - 5 .
\]
\end{exercise}

\begin{exercise}[$\star\star$]\label{exo:g9:algebra:9}
Je pense à un nombre, je le multiplie par $3$, j'ajoute $7$, et j'obtiens
le même résultat que si je l'avais multiplié par $5$ et soustrait $9$.
Quel est mon nombre~? (Mettre en place une équation et la résoudre.)
\end{exercise}

\begin{exercise}[$\star\star\star$]\label{exo:g9:algebra:10}
Soit $n$ un nombre entier.
\begin{enumerate}
  \item Développer $(n + 1)(n - 1)$ et en déduire que le produit des
        deux voisins de $n$ est toujours $n^2 - 1$ (par ex.\
        $24 \times 26 = 25^2 - 1 = 624$).
  \item Montrer que la somme de trois entiers consécutifs est toujours
        un multiple de $3$.
\end{enumerate}
\end{exercise}
